% Proposed B.120. Owned fragment only; no assembled source/register mutation.
\subsection*{The intermediate boundary amalgam and its four remaining relations}

\statement{B.120}{Proved}{The actual boundary amalgam has a proper full completion, an exact discrete free-six cone centralizer, and a fully embedded carrier for the four remaining prefix relations}
\label{stmt:B120}
Use the actual regional marking of B.112, B.114 and B.118. To
distinguish the adapted cone from the selected upper-gap cone, write
\[
 \bar x=CxC,\quad B=CbC,\quad z=Ch,\quad f=hk^{-1},\quad
 K=F(a,B,d,e,f),\quad c=aCa^{-1}.
\]
The actual physical neck basis and its prefixes are
\begin{equation}\label{eq:B120-prefixes}
 \begin{gathered}
 (O_1,O_2,O_3,O_4,O_5)
   =(aBa^{-1},a,fde d^{-1}f^{-1},d,d^{-1}fd),\\
 P_i=O_1\cdots O_i,\qquad P=P_5=aBfde,\qquad T=cP.
 \end{gathered}
\end{equation}
In particular $P_1,\ldots,P_5$ are a free basis of the actual $K$.
Let $R$ be the whole regional group, including every annular node
relation and the retained zero-loop action. By B.118 its actual
boundary subgroup is
\[
 H=\langle c,k,T\rangle=C_2(c)*\mathbb Z^2(k,T),\qquad
 K\cap H=\langle P\rangle,
\]
and $H$ and $K$ have their full induced topology in $R$, with
$\widehat K\cap\widehat H=\overline{\langle P\rangle}$ there.
Put
\begin{equation}\label{eq:B120-corner-amalgam}
 \begin{split}
 W_{\partial}&=\langle D,u,k,j\mid D^2=j^2=1,
              \ [k,D]=[k,u]=1,\ juj=u^{-1}\rangle,\\
 M&=R*_{H}W_{\partial},
 \end{split}
\end{equation}
where the corner edge map is $c\mapsto jDj$, $T\mapsto uD$,
and $k\mapsto k$. Both are the actual embedded edge maps of
B.118. Then the full completion is a proper profinite amalgam:
\begin{equation}\label{eq:B120-full-completion}
 \widehat M\cong
       \widehat R\amalg_{\widehat H}\widehat{W_{\partial}}.
\end{equation}
Here both vertex completions, and hence the indicated edge
completion, inject.

Let $L\cong F_5$ be the actual selected regional source in $R$,
so that $C_R(c)=L\times\langle c\rangle$ by B.112 after the
specified conjugation by $a$. With $k'=jkj$ one has the exact
\emph{discrete} formulas
\begin{equation}\label{eq:B120-discrete-centralizers}
 \begin{aligned}
 C_M(c)&=(L*\langle k'\rangle)\times\langle c\rangle
                      \cong F_6\times C_2,\\
 C_M(D)&=(jLj*\langle k\rangle)\times\langle D\rangle.
 \end{aligned}
\end{equation}
The full actual mirror extension $X$ of the annular one-neck group
is precisely
\begin{equation}\label{eq:B120-final-quotient}
 X=M/\langle\!\langle(jP_i)^2:\ 1\leq i\leq4\rangle\!\rangle.
\end{equation}
The four relations lie in the genuinely embedded subgroup
\begin{equation}\label{eq:B120-carrier}
 J_0=\langle K,j\rangle
    =F(w_1,w_2,w_3,w_4)*C_2(j)*C_2(jP_5),\qquad w_i=jP_i.
\end{equation}
This carrier has its full induced topology in $M$, so its full
completion also embeds in $\widehat M$.
The six words identified below are actual source prelifts in the
free factor of $C_M(D)$; their images under
\eqref{eq:B120-final-quotient} are the six original based source
words. Neither a centralizer formula in $\widehat M$ nor survival
of these source generators after the four relations is asserted.

\begin{proof}
\emph{B.120.1. The exact intermediate presentation.}
The two actual edge embeddings from B.118 and ordinary amalgam
normal form make \eqref{eq:B120-corner-amalgam} a proper discrete
amalgam. Eliminate $D=j c j$ and $u=T D$. The relation
$[k,D]=1$ becomes $[jcj,k]=1$, while $[k,u]=1$ follows from
this equation and the already established $[k,T]=1$. Finally,
using $P=cT$, the equation $juj=u^{-1}$ is equivalent to
$jPj=P^{-1}$. Indeed
\[
 P=jDj\,uD,
 \qquad jPj=P^{-1}
 \quad\Longleftrightarrow\quad juj=u^{-1},
\]
by cancellation with $D^2=j^2=1$. Consequently
\begin{equation}\label{eq:B120-intermediate-presentation}
 M=\langle R,j\mid j^2=1,\ [jcj,k]=1,\ jP_5j=P_5^{-1}\rangle.
\end{equation}
This is a reversible change of presentation.

The actual reflected neck automorphism satisfies
$\eta(P_i)=P_i^{-1}$ on the full actual prefix basis. Before
imposing the two cross commutators, folding the two regional
vertices gives $R*_{K}(K\rtimes_\eta C_2(j))$. The first cross
commutator is $[jcj,k]$; its $j$-conjugate is the other actual
cross commutator. Thus \eqref{eq:B120-intermediate-presentation}
retains the complete regional presentation, the cross relations,
and the fifth prefix relation. Adding the first four prefix
relations gives exactly \eqref{eq:B120-final-quotient}. This
argument uses the actual based reflection of B.114, not a separate
inversion of individual neck generators. The original annular
one-neck group is the kernel of the mirror-parity map $X\to C_2$
which is zero on $R$ and one on $j$. This parity is distinct from
the original ordinary-cover character.

\emph{B.120.2. A corner companion that extends over the whole region.}
Let $W_0$ be the mirror-even subgroup of $W_{\partial}$. Its exact
index-two rewriting is the path graph-product presentation
\begin{equation}\label{eq:B120-even-corner}
 W_0=\langle D,k,u,k',c\mid D^2=c^2=1,
 [D,k]=[k,u]=[u,k']=[k',c]=1\rangle,
 \quad k'=jkj.
\end{equation}
There is a retraction $\rho:W_0\to H$ defined by
\begin{equation}\label{eq:B120-corner-retraction}
 \rho(D)=\rho(k')=1,\quad
 \rho(u)=T,\quad\rho(c)=c,\quad\rho(k)=k.
\end{equation}
Every relation in \eqref{eq:B120-even-corner} is respected;
moreover $\rho(T=uD)=T$. On $H$ the companion
$\nu=(\rho\circ\operatorname{Ad}(j))|_H$ is
\begin{equation}\label{eq:B120-nu}
 \nu(c)=\nu(k)=1,\qquad
 \nu(T)=T^{-1}c=P^{-1},\qquad \nu(P)=P^{-1},\qquad \nu^3=\nu.
\end{equation}
For example $jTj=u^{-1}c$, giving the displayed value of $\nu(T)$.

This particular companion extends to the actual regional group.
Use the presentation character
\[
 \lambda:R\longrightarrow\mathbb Z,\qquad
 \lambda(h)=1,\quad
 \lambda(C)=\lambda(k)=\lambda(a)=\lambda(b)=\lambda(d)=\lambda(e)=0.
\]
All relations in B.112's finite HNN presentation respect this
character. It is the branch-height character used in that proof,
not the geometric annular-base character. Since $f=hk^{-1}$,
one has $\lambda(P)=\lambda(T)=1$ and $\lambda(c)=0$. Therefore
\begin{equation}\label{eq:B120-extension}
 v:R\longrightarrow R,\qquad v(g)=P^{-\lambda(g)}
\end{equation}
is an actual homomorphism with image in $H$, satisfying
$v|_H=\nu$ and $v^3=v$. No arbitrary retraction $R\to H$ is
being assumed.

\emph{B.120.3. Cofinal finite kernels with the same actual edge kernel.}
Let $A_R\lhd R$ and $A_W\lhd W_{\partial}$ be arbitrary prescribed
finite-index normal subgroups. The full induced topology of $H$
in $R$ supplies $V_R\lhd R$ of finite index such that
\[
 V_R\leq A_R,\qquad V_R\cap H\leq A_W\cap H.
\]
Write $\psi:R\to R/V_R$ and $N=V_R\cap H$. Form the finite
product homomorphism
\begin{equation}\label{eq:B120-triple-map}
 \Phi=(\psi,\psi v,\psi v^2):R\longrightarrow(R/V_R)^3.
\end{equation}
Its restriction to $H$ has exactly the kernel
\begin{equation}\label{eq:B120-edge-kernel}
 L_0=N\cap\nu^{-1}(N)\cap\nu^{-2}(N).
\end{equation}
Thus $L_0\lhd H$ has finite index, $L_0\leq A_W\cap H$, and
$\ker\Phi\leq A_R$. Also $\nu(L_0)\leq L_0$: the three
conditions for $\nu(x)$ follow from the three conditions for $x$
and $\nu^3=\nu$.

The genuine paired homomorphism on $W_0$ is
\[
 b\longmapsto
       (\rho(b)\bmod L_0,\ \rho(jbj)\bmod L_0).
\]
Conjugation by $j$ interchanges its coordinates, so it extends to
\begin{equation}\label{eq:B120-paired-map}
 \Psi:W_{\partial}=W_0\rtimes\langle j\rangle
       \longrightarrow (H/L_0\times H/L_0)\rtimes C_2,
 \qquad j\longmapsto\text{coordinate swap}.
\end{equation}
On $H$ it is $x\mapsto(x\bmod L_0,\nu(x)\bmod L_0)$.
Its restriction kernel is exactly $L_0$, using the proved
$\nu$-stability. Put
\[
 V_W=\ker\Psi\cap A_W.
\]
Then $V_W\lhd W_{\partial}$ has finite index, refines $A_W$,
and still satisfies $V_W\cap H=L_0$, because $L_0\leq A_W$.
Thus the finite vertex quotients
\[
 F_R=R/\ker\Phi,\qquad F_W=W_{\partial}/V_W
\]
contain precisely the same embedded edge quotient $H/L_0$.

We recall explicitly the finite-target step. For two finite groups
with a common embedded subgroup, let each act freely on a finite
set of cardinality divisible by both group orders. Their restricted
edge actions are unions of equally many regular edge actions.
Conjugating one action aligns those restrictions and gives a
finite permutation image of their amalgam which is injective on
both finite vertices. Applied to
\[
 F_R*_{H/L_0}F_W,
\]
this produces a finite target of $M$ refining both prescribed
vertex kernels. In particular $R$ and $W_{\partial}$ have their
full induced topology in $M$. These finite-group amalgams are
also residually finite: the just-constructed finite quotient has
a finite-index free kernel, by its free action on the ordinary
amalgam tree; characteristic finite-index subgroups of that
finitely generated free kernel separate its elements.

Every finite target of $M$ factors through one of the displayed
finite-group amalgams: apply the construction to the two
restriction kernels of that target. The universal property for
finite targets, followed by inverse limits, now gives
\eqref{eq:B120-full-completion}. The vertex injections follow
from the explicit cofinal refinement, so this profinite pushout
is proper. The proof does not assert residual finiteness of $M$
or separability of arbitrary subgroup products.

\emph{B.120.4. The discrete cone centralizer.}
Ordinary free-product normal form in
$H=C_2(c)*\mathbb Z^2(k,T)$ shows that every involution of $H$
is $H$-conjugate to $c$ and that $C_H(c)=\langle c\rangle$.
On the regional side B.112 gives
$C_R(c)=L\times\langle c\rangle$. On the corner side, write
\[
 W_{\partial}=
   ((C_2(D)*\mathbb Z(u))\times\mathbb Z(k))
       *_{\mathbb Z(u)}(\mathbb Z(u)\rtimes C_2(j)).
\]
The edge is torsion free. The involution $D$ fixes just its base
vertex in this ordinary tree, and its centralizer there is
$\langle D\rangle\times\langle k\rangle$. Conjugating by $j$
gives $C_{W_{\partial}}(c)=\langle c\rangle\times\langle k'\rangle$.

In the ordinary tree of $M$, the $c$-fixed subtree contains the
base edge. At either endpoint, consider any incident $c$-fixed
edge represented by $gH$ in that vertex group. Then
$g^{-1}cg\in H$ is an involution. Its unique $H$-conjugacy class
allows one to choose $t\in H$ with
$g^{-1}cg=tct^{-1}$. The representative $gt$ centralizes $c$
and still represents the same edge. Thus the vertex centralizer
is transitive on its incident fixed edges. Applying this adjustment
along each finite path in the fixed subtree makes the centralizer
quotient exactly the base edge. Equivalently, reducing an amalgam
word with these successive adjustments gives
\[
 C_M(c)=C_R(c)*_{\langle c\rangle}C_{W_{\partial}}(c).
\]
The common involution is central in both factors, yielding the
first formula in \eqref{eq:B120-discrete-centralizers}.
Its conjugate by $j$ gives the second. This is an ordinary-tree
argument. Properness of \eqref{eq:B120-full-completion} alone
does not transfer it to a completed centralizer.

\emph{B.120.5. Literal source prelifts and the final source map.}
Here are the raw positive regional source words, before their
triangular simplification in B.112. Set
\[
 E=fde d^{-1}f^{-1},\qquad
 \Theta=CbCb\,CECb^{-1}Cb^{-1},\qquad \delta_U=U\bar U.
\]
In the order of the four contacts followed by the zero source loop,
put
\begin{equation}\label{eq:B120-raw-source}
 \begin{gathered}
 q_1=\bar b b,\qquad q_2=\Theta^2,\\
 q_3=Cb[d,\bar d]b^{-1}C\,\delta_d,\qquad
 q_4=Cb[a,\bar a]b^{-1}C\,\delta_a,\qquad q_5=h,\\
 s_i=a q_i a^{-1}\quad(1\leq i\leq5).
 \end{gathered}
\end{equation}
The $s_i$ are a free basis of the actual $L$. This is the explicit
source-basis comparison of B.112, transported by the selected
upper-gap connector; it is not a rank argument. Consequently
\begin{equation}\label{eq:B120-six-prelifts}
 js_1j,\ js_2j,\ js_3j,\ js_4j,\ js_5j,\ k
\end{equation}
are a free basis of the ordinary factor of $C_M(D)$.

Under the quotient \eqref{eq:B120-final-quotient}, and only there,
the five prefix relations imply $jxj=\eta(x)$ for every $x\in K$.
The plus regional tuple is then literally the $j$-image of the
minus tuple, with the transported selected connector intact.
Thus \eqref{eq:B120-six-prelifts} maps to the five actual outer
regional source words followed by the actual inner good source
loop. This also follows by direct free-word substitution in the
unchanged actual nine-generator package: the independent checker
\path{work_A/code/check_intermediate_source_prelifts25.py}
checks all nine inverse mirror substitutions, all six raw source
comparisons, and the four remaining prefix relators. The full
words and input hash are retained in
\path{work_A/intermediate_source_prelifts25.json}; genuine output
records prelift lengths $8,28,18,18,6,1$ and final raw lengths
$16,64,102,42,24,1$. The checker applies $\eta$ only in the
explicit final quotient map and imports no earlier producer.
No equality between $jKj$ and $\eta(K)$ is used inside $M$.

\emph{B.120.6. An embedded carrier for exactly four square relations.}
Inside $W_{\partial}$ put $Q_0=\langle P,j\rangle$.
It is the infinite dihedral group on these marked generators.
Indeed, $jPj=P^{-1}$ holds, and the map
\[
 W_{\partial}\longrightarrow Q_0,\qquad
 D,k\longmapsto1,\quad u\longmapsto P,\quad j\longmapsto j
\]
is a retraction: the actual word $P=jDj\,uD$ maps to $P$.
Equivalently the same assignments into the abstract infinite
dihedral group give a left inverse to its marked presentation
map. Hence $Q_0$ also has its full induced topology in the corner.
The mirror-even subgroup of $Q_0$ is exactly $\langle P\rangle$,
whereas $H$ is mirror even. It follows that
\begin{equation}\label{eq:B120-corner-face}
 Q_0\cap H=\langle P\rangle,\qquad
 \widehat{Q_0}\cap\widehat H=\overline{\langle P\rangle}
       \quad\text{inside }\widehat{W_{\partial}}.
\end{equation}
For the completed equality use the continuous parity character:
its kernel on the embedded $\widehat{Q_0}$ is precisely the full
procyclic closure of $\langle P\rangle$, already contained in
$\widehat H$. No exchange of completion and intersection is used.

Together with $K\cap H=\langle P\rangle$ on the regional side,
ordinary amalgam normal form now embeds
\begin{equation}\label{eq:B120-carrier-amalgam}
 K*_{\langle P\rangle}Q_0\longrightarrow M.
\end{equation}
Every alternating syllable outside $\langle P\rangle$ is outside
$H$ in its corresponding vertex, so every reduced word remains
reduced in $M$. This gives the complete discrete injection proof.
Because the actual prefixes are a free basis of $K$, its image is
\[
 F(P_1,P_2,P_3,P_4)*D_{\infty}(P_5,j).
\]
Multiplying each of the first four free generators on the left
by $j$ is an invertible free-product Nielsen change, fixing the
dihedral factor. This proves \eqref{eq:B120-carrier}.

Each $P_i$ for $i<5$ lies outside $H$, by the proved exact face
$K\cap H=\langle P_5\rangle$, and $j$ lies outside $H$ by
mirror parity. Thus $(jP_i)^2$ is cyclically reduced of length
four in the actual amalgam, and $jP_i$ has tree translation length
two. The quotient of $J_0$ by these four squares is precisely
\[
 Q=C_2(w_1)*C_2(w_2)*C_2(w_3)*C_2(w_4)*C_2(j)*C_2(jP_5),
\]
the actual mirror neck group. Accordingly $X$ is the categorical
pushout $M*_{J_0}Q$, with a \emph{surjective} map $J_0\to Q$.
It is not an amalgam with two injective edge maps. The completed
face identities on the two vertices do not by themselves establish
the carrier's full topology; we prove that next by an additional
actual quotient and finite-target construction.

\emph{B.120.7. Full topology of the carrier by an actual ordinary retraction.}
The ordinary regional retraction of B.112--B.114 is, in the
present neck coordinates, an actual homomorphism
\begin{equation}\label{eq:B120-ordinary-retraction}
 \begin{gathered}
 r:R^0\longrightarrow K,\qquad r|_K=1_K,\\
 r(CaC)=r(b)=r(CdC)=r(CeC)=1,\qquad r(CfC)=r(h)=f.
 \end{gathered}
\end{equation}
Here $R^0$ denotes the original ordinary index-two subgroup,
not the mirror-even subgroup. These values are the complete
regional detector already checked on both generic action
directions and all twenty ordinary node relations in B.112.
In particular $r(k)=r(f^{-1}h)=1$. Since $ChC=h$,
one also has $r(CkC)=r((CfC)^{-1}h)=1$. Finally,
\begin{equation}\label{eq:B120-conjugate-prefix-retraction}
 r(P)=P,\qquad r(CPC)=f.
\end{equation}
Indeed $P=aBfde$ and the conjugate of each of $a,B,d,e$
has trivial image under $r$, whereas $CfC$ maps to $f$.

Kill the normal closure of $k$ in the entire intermediate group.
Writing $\bar R=R/\langle\!\langle k\rangle\!\rangle_R$, the
result is exactly
\begin{equation}\label{eq:B120-killed-base}
 \bar M=\langle\bar R,j\mid j^2=1,\ jPj=P^{-1}\rangle.
\end{equation}
The cross commutator disappears and no other relation is added.
The character $\varepsilon(C)=1$, $\varepsilon(j)=0$, zero on
the ordinary regional generators, still maps $\bar M$ onto
$C_2$. Its index-two kernel has the exact presentation
\begin{equation}\label{eq:B120-ordinary-killed-presentation}
 \begin{gathered}
 \bar M^0=
 \langle\bar R^0,j_0,j_1\mid
       j_0^2=j_1^2=1,\ j_0Pj_0=P^{-1},\ j_1Qj_1=Q^{-1}\rangle,\\
 j_0=j,\qquad j_1=CjC,\qquad Q=CPC.
 \end{gathered}
\end{equation}
This is the two-coset rewriting with representatives $1,C$:
the two copies of the inversion relation are exactly the two
displayed equations. Because $r$ kills $k$ and $CkC$, it factors
through $\bar R^0$. To see that this accounts for the whole
quotient, the normal closure of $k$ in $R$ is the normal closure
inside $R^0$ of these two elements, using $R=R^0\rtimes\langle C\rangle$.

By \eqref{eq:B120-conjugate-prefix-retraction} there is now a
homomorphism
\begin{equation}\label{eq:B120-carrier-target}
 \bar M^0\longrightarrow
 \mathscr L=J_0*_{\langle f\rangle}
             (\langle f\rangle\rtimes C_2(j_1)),
 \qquad j_1 f j_1=f^{-1},
\end{equation}
whose restriction to $\bar R^0$ is $r$ and which sends $j_0$
to the actual $j$ of $J_0$. The map on the natural copy of $J_0$
is the identity. In particular this construction also verifies
directly that the carrier survives the quotient killing $k$.
The edge here is genuinely infinite cyclic: $f$ has infinite
order in the actual free $K\leq J_0$.

Let $\chi:J_0\to F$ be any finite target, replacing $F$ by its
image if necessary, and let $m$ be the order of $\chi(f)$.
Then \eqref{eq:B120-carrier-target} maps into the finite-vertex
amalgam
\[
 F*_{\langle\chi(f)\rangle}(C_m\rtimes C_2),
\]
where the involution acts by inversion and the cyclic factor is
identified with $\langle\chi(f)\rangle$. This includes $m=1$.
The finite permutation construction in B.120.3 supplies a finite
quotient injective on $F$. Composing gives a finite quotient of
$\bar M^0$ whose restriction kernel on $J_0$ is exactly
$\ker\chi$. Core its kernel through the actual index-two
extension $\bar M^0\leq\bar M$, and pull the resulting finite
quotient back to $M$. The restriction kernel on $J_0$ can only
decrease, and so is contained in $\ker\chi$.
Every finite-index normal subgroup of $J_0$ is therefore refined
by a finite-index normal subgroup of $M$. This proves the full
induced topology and the claimed injection
\begin{equation}\label{eq:B120-carrier-full-topology}
 \widehat{J_0}\hookrightarrow\widehat M.
\end{equation}
This additional argument uses the actual ordinary regional
retraction, not a general profinite subamalgam assertion.

\emph{B.120.8. The remaining filling issue.}
The intermediate group has the exact discrete source centralizer
and the fully embedded completed vertices proved above. The four
remaining actual square relations can still create new centralizing
elements or impose relations on the image of the six source
prelifts. No relative small-cancellation, congruence-extension,
or completed-centralizer conclusion is inferred from their length
or from the carrier's free-product decomposition. The selected
source comparison in the actual quotient remains a separate gate.
\end{proof}
